\documentclass[a4paper,10pt,twocolumn,landscape]{article}
	\usepackage[ansinew]{inputenc}
	\usepackage{ae}
	\usepackage[brazil]{babel}
	\usepackage{xcolor}
	\usepackage{calc}
	\usepackage{wallpaper}
	\usepackage[margin=2cm]{geometry}
	\usepackage{SIunits}
	\usepackage{amsmath}
	\usepackage{listings}


% Configurações personalizadas do código LaTeX.	
\lstnewenvironment{LaTeXcode}{
	\setlength{\abovecaptionskip}{0pt}
	\lstset{language=[LaTeX]TeX}
	\lstset{emph={&},emphstyle={\textcolor{red}}}
	\lstset{%
		basicstyle=\footnotesize\ttfamily,  % Global
		keywordstyle=\color{blue}\bfseries, % Comandos
		identifierstyle=,                   % Texto
		stringstyle=,                       % Strings 
		commentstyle=\color{gray},          % Comentários
		showstringspaces=false,             % Espaços
		rulecolor=\color{gray},             % Linha da caixa
	}
}% Abrindo o ambiente.
{}% Fechando o ambiente.
	
\newcommand{\pkg}[1]{{\normalfont\sffamily\color{orange}#1}}
\newcommand{\env}[1]{{\normalfont\sffamily\color{green!50!black}#1}}

	% Papel de parede
	\CenterWallPaper{1}{../figuras/leao-pensador-wattermark.png}

	% Configuração extra: espaço entre as colunas.
	\setlength{\columnsep}{2cm}

	\newcommand\vetor[1]{\mathbf{#1}}

\begin{document}

\begin{center}\includegraphics[width=0.8\columnwidth]{../figuras/LogotipoCursoLaTeX_v2}\end{center}

	O pacote \pkg{amsmath} define os ambientes \env{align} e \env{align*} (o asterísco faz parte do nome do ambiente), que lhe permitem alinhar várias expressões verticalmente. Por exemplo, a expressão
	\begin{align*}
	I_a &= I_b + M \left(\mathbf{R} \times \mathbf{n}\right)^2\\
		 &= I_b + M R^2 \sin^2\theta
	\end{align*}
é produzida da seguinte forma:
	\begin{LaTeXcode}
\begin{align*}
I_a &= I_b + M \left(\mathbf{R} \times \mathbf{n}\right)^2 \\
    &= I_b + M R^2 \sin^2\theta
\end{align*}
	\end{LaTeXcode}

	A novidade aqui são as instruções \verb|&| e \verb|\\|. A primeira marca a posição do alinhamento: imediatamente antes do sinal de igual. A segunda, como você já deve ter imaginado, provoca a mudança de linha.

	A diferença entre \env{align} e \env{align*} é que o primeiro insere automaticamente um número a cada linha da expressão, independentemente de você ter ou não atribuido um nome (\verb|\label|) a ela. Inversamente, \env{align*} \emph{não} insere números.

	Com essas informações, produza a seguinte expressão:
	\begin{align}
		\int_{\vetor r_0}^{\vetor r_1} \vetor\nabla \varphi \cdot d\vetor l 
				&= \varphi(\vetor{r}_2) - \varphi(\vetor{r}_1)\label{eq:calc} \\
		\int_S \vetor\nabla \times \vetor F \cdot \vetor n\, d\vetor a
				&= \oint_C \vetor F \cdot d\vetor l\label{eq:Stokes} \\
		\int_V \vetor\nabla \cdot \vetor F \, dv
				&= \oint_S \vetor F \cdot \vetor n\, da.\label{eq:divergente}
	\end{align}

	A equação \ref{eq:calc} é o teorema fundamental do cálculo, \eqref{eq:Stokes} é o teorema de Stokes e \eqref{eq:divergente}, o do divergente.

	Atenção para os símbolos em negrito, que representam grandezas vetoriais, produzidos com o comando \verb|\mathbf|. Por falar nisso, responda: por que não utilizar \verb|\textbf|? 

	Você já sabe que no \LaTeX\ devemos sempre procurar definir a função dos elementos do texto (ou das expressões matemáticas). No entanto, \verb|\mathbf| define a forma (ou leiaute), à imagem de \verb|\textit|. Considerando o que vimos na atividade 8 e na sua discussão com seus colegas \TeX starters na tarefa conjunta, como você poderia resolver este impasse? Você não precisa \emph{fazer} isso, mas apenas explicar.

	Responda essas perguntas e produza a expressão matemática acima num arquivo \LaTeX\ e envie-o para o professor/tutor.

\end{document}
